% section: 漸化式の分類

\section{漸化式の分類}
具体的な解法について触れる前に,解法の分類を行おう.
まず漸化式は解ける漸化式と解けない漸化式に分類される.
そして漸化式は「決まった手順で解ける漸化式」と「特定の性質から解ける漸化式」に分類される.
必ず解ける漸化式はさらに分類され,各分類ごとに様々な解法がある.
ここで言う「解ける」というのは,本書内での定義であり,厳密な定義は存在しない.
一応の基準として高校数学程度の範囲内で既知の関数に帰着するか,有限回の演算または操作により一般項を得られるような解法が存在するものを意味している.

本書では,漸化式の基本的なパターンを説明し,次にそこからいかにしてそのパターンに帰着させるかを議論する.
それが応用パターンの節である.
